% Untracked draft: do not commit
\usepackage{amssymb}
\usepackage{pifont}

\subsection*{I.B Semantic lowering and provenance}

Raw graph inventory is not itself a claim; meaning enters the pipeline only through deterministic semantic-lowering rules applied to the recorded nodes, and this layer constitutes the core contribution of the current system. Each rule recognizes one specific structural construction and derives one semantic value from it. An adjacency-fed message stage is identified when a matrix-multiply operator (\texttt{aten.matmul.default}, \texttt{aten.mm.default}, \texttt{aten.bmm.default}, or \texttt{aten.mv.default}) consumes exactly one reference to the declared adjacency state, or to an existing alias of it, together with exactly one reference to a previous message state; unrelated matrix multiplications are never counted toward depth. The message-passing depth is then obtained by tracing the chain of such stages backward from the declared \texttt{message\_state} output and taking its length, terminating at the first non-matching or ambiguous node rather than speculating past it. A hinge-like XC construction is recognized when an ancestor of the \texttt{xc\_energy} output applies a recognized nonsmooth activation such as \texttt{aten.relu.default}, \texttt{aten.leaky\_relu.default}, \texttt{aten.clamp\_min.default}, \texttt{aten.maximum.default}, \texttt{aten.abs.default}, or \texttt{aten.hardtanh.default}, in which case the XC form is classified as a hinge and the architecture is taken to contain the nonsmooth construction required for discontinuity support; conversely, if no hinge ancestor exists but an ancestor applies a smooth activation (\texttt{sigmoid}, \texttt{softplus}, \texttt{tanh}, \texttt{silu}, or \texttt{gelu}), the XC form is classified as smooth, a classification that explicitly does not support the required discontinuity. Operator constructions are classified at their root: \texttt{zeros} yields a structurally self-adjoint zero operator, \texttt{eye} an identity, and a root of the additive form $B + B^{\mathsf{T}}$, in which one operand is a direct transpose, permute, or \texttt{T} application to the same base tensor as the other operand, yields a symmetrized construction whose self-adjointness follows structurally, whereas a root that is merely a placeholder attribute or fetched parameter yields an unconstrained parameter with no self-adjointness guarantee.

Every derivation records its full provenance, because a semantic assertion is only auditable if it can be traced back to the exact nodes that justify it. The derivation record therefore contains the derived value, the rule identifier and version that produced it, the root node at which the pattern was matched, the supporting graph nodes that constitute the evidence, and the observed operations actually present in the inventory. These records are not annotations appended after analysis but constitutive of the claim itself, so that any semantic assertion in the final certificate can be traced through the IR to specific nodes of the hashed artifact. Matching is exact on lower-cased targets: unlisted spellings, custom operators, and future Torch variants are never accepted on the basis of name similarity, structural plausibility, or reviewer intuition about what the model was intended to compute.

The governing principle of this layer is fail-closed behavior, meaning that compositions matching none of the supported patterns are emitted as \texttt{unsupported} rather than guessed. A graph whose operator root matches neither a recognized construction nor a placeholder produces no self-adjointness claim whatsoever; a message chain interrupted by an ambiguous or non-adjacency-fed multiplication contributes depth only up to the interruption; an XC head built from unrecognized operations yields no form classification at all. The topology derivation is similarly strict, requiring a non-empty square boolean or integer structural matrix in which every truthy entry becomes one directed edge under the declared convention, defaulting to target-source ordering unless source-target is explicitly declared. The system therefore certifies precisely those structures it can name and declines everything else rather than extending partial credit to near-matches; the cost of this conservatism is rejection of unusual architectures, and the benefit is that every issued claim corresponds to an explicit, versioned, auditable rule.

\subsection*{II.C Structural IR}

Lean does not receive the Torch graph, and nothing in the proof layer depends on the semantics of any deep-learning framework. The proof layer receives instead a compact intermediate representation, conceptually
\[
\mathcal{I} = (E,\, d,\, Q,\, x,\, o),
\]
where $E$ is the finite directed topology derived from the adjacency state under the declared convention, $d$ is the message-passing depth from the backward trace, $Q$ is the set of required source/target couplings whose coverage is to be proved, $x \in \{\text{hinge}, \text{smooth}\}$ is the XC form, and $o \in \{\text{zero}, \text{identity}, \text{symmetrized}, \text{unconstrained}\}$ is the operator construction. Each component carries its derivation record and node-level provenance from Section I.B, so the IR is simultaneously the object reasoned about and the index into the evidence. The representation is deliberately small and entirely finite, so that every property subsequently claimed about it reduces to a decidable obligation over explicit data, such as coverage of finitely many directed pairs by finitely many stages, rather than to reasoning about floating-point execution or framework internals.

\subsection*{II.D Translation validation}

Because the lowering from inventory to IR is trusted Python, the pipeline validates it by deterministic recomputation rather than by formal proof. From the raw inventory hash and the reviewed input constraints, a second independent pass recomputes the expected semantic derivations for topology, message depth, XC form, and operator construction, together with the output roots, the adjacency state and its aliases, and the declared conventions, and then compares these expected results against the stored IR; any disagreement between inventory, derivation, and IR causes the translation to be rejected outright. This duplicate evaluation catches implementation drift, version skew between analyzer components, and accidental mutation of stored results, at the cost of doubling a cheap computation.

Two distinctions must be kept explicit when assessing what this validation establishes. First, hash integrity is not semantic validity: the SHA-256 binding establishes which artifact was analyzed, while the semantic rules determine what that artifact is allowed to assert in the IR, and neither property subsumes the other. Second, the Torch-to-IR semantic lowering remains part of the trusted base; Lean proves consequences of the encoded finite structure, not the correctness of the Python program that produced the encoding. Supporting a new operator consequently requires updating the written specification and the analyzer version, adding both a positive mapping test and a near-name or custom-operator rejection test, retaining node-level provenance and hash binding throughout, and confirming that the existing generic Lean predicate expresses the intended structural consequence.

\subsection*{II.E Lean obligation generation}

Given $\mathcal{I}$ and a certification requirement, the system generates a fixed Lean theorem statement covering coupling coverage over $E$ and $Q$, discontinuity support conditional on $x = \text{hinge}$, and self-adjointness conditional on $o \in \{\text{zero}, \text{identity}, \text{symmetrized}\}$. The theorem statement is generated deterministically; proof-search components may propose proofs but cannot redefine the verification target. A proposed proof body is accepted only if it compiles against the fixed statement under the Lean kernel, so the soundness of the resulting certificate reduces to kernel checking plus the trusted lowering discussed above, not to the behavior, completeness, or trustworthiness of any search procedure. Human-approved claims occupy a strictly intermediate position in this scheme: they are input to proof generation and can never substitute for an accepted proof body, which is why even fully confirmed specifications receive a permanent certificate kind distinct from artifact-grounded certification.

\subsection*{II.F Certificate and scope}

The issued certificate binds, in a single attested object, the artifact digest $H_{\text{artifact}}$, the IR digest $H_{\text{IR}}$, the rule identifiers and versions used in the lowering, the generated obligations, and the Lean acceptance results. Any change to the artifact, the specification, the rules, or the obligations invalidates the binding, so the certificate is meaningful only relative to the exact byte content it names.

The scope of the resulting guarantee is exactly structural:
\begin{itemize}
  \item structural certification: \checkmark
  \item trained weights: \ding{55}
  \item numerical accuracy: \ding{55}
  \item convergence: \ding{55}
  \item experiment agreement: \ding{55}
\end{itemize}
